% Complete all-M proof, integrated in ../paper.tex. Labels AM refer to this section.
\section{Completion for every cluster size on branches 1--3}\label{sec:allm}

The certificates below remove both restrictions of the original finite-range result:
the missing spectral hypothesis for intermediate $M$ and the unbounded range $M>10^5$.
All stability statements concern the invariant mass hyperplane through the constant equilibrium.
The term ``branch 1'' identifies the continuation from the first limiting Hopf root; it does not
assert stability for every parameter below that point.

\textbf{Theorem D (complete local criticality on branches 1--3).}
For each $j=1,2,3$, the Hopf curve continuing the $j$th root of the Pego--Velázquez
limit characteristic function is a graph over $M\in(M_j^*,\infty)$, with
\[
 14<M_1^*<15,\qquad43<M_2^*<44,\qquad86<M_3^*<87.
\]
It has exactly one point at every integer $M\ge M_j$, where $(M_1,M_2,M_3)=(15,44,87)$,
and none below this onset. At every such point the imaginary pair is simple, crosses
transversally towards the right as $\kappa$ increases, and is the only spectrum on the
imaginary axis after restriction to $H$. The first Lyapunov coefficient is negative
except at $(j,M)=(1,15)$, where it is positive. Consequently:
\begin{itemize}
\item Branch 1 has orbitally asymptotically stable small periodic orbits for
      $\kappa>\kappa_1(M)$ close to the bifurcation, for every $M\ge16$.
      At $M=15$ its small periodic orbits are subcritical and unstable.
\item Branches 2 and 3 are supercritical on the two-dimensional centre manifold,
      for every $M\ge44$ and $M\ge87$, respectively. Their small periodic orbits
      are unstable in $H$ because a transverse unstable eigenvalue is present.
\item For $3\le M\le14$ the equilibrium is linearly asymptotically stable on $H$
      for every $K>0$; there is no Hopf bifurcation.
\end{itemize}
The uniqueness here is for the specified connected Hopf curves, not a classification
of all possible disconnected components of the Hopf set.

\subsection{AM1. A characteristic function regular at infinity}
Put $\varepsilon=M^{-1/2}$, $K=\kappa\varepsilon$, $A=1+K$, $\varphi=1+\varepsilon^2z$.
Pego--Velázquez's characteristic determinant has the exact form
\[
 F=-P_1+\varphi^MP_2+A^{-M}R_1+(A\varphi)^{-M}R_2,
\]
where, writing $S_0=(\varphi-1)(A\varphi-1)(A\varphi^2-1)$,
\begin{align*}
 P_1&=K[(A\varphi-1)^2\varphi^2-(\varphi-1)^2]+(KM+1)S_0,\\
 P_2&=\varphi(A\varphi-1)^3-A^2(\varphi-1)^2+MA^2(A\varphi-1)(\varphi-1)^2,\\
 R_1&=-A^2[(A\varphi-1)^2\varphi^2-(\varphi-1)^2],\\
 R_2&=A(\varphi-1)^3+A^2\varphi(A\varphi-1)^2[M(\varphi-1)+\varphi].
\end{align*}
Except at the real exceptional points $1,1/A,\pm A^{-1/2}$, zeros correspond to eigenvalues
$\lambda=(A-\varphi^{-1})(\varphi-1)$. A simple zero gives a simple eigenvalue when
$\varphi f(\varphi)=A\varphi^M+\varphi-1\ne0$, as checked below; this supplies the
extra condition in the corrected application of PV Lemma 5.4. The eigenvalue zero of $B$
is simple and is removed on $H$ (PV Lemma 3.1).

Define the analytic functions
\[
 \mathsf L(u)=\frac{\log(1+u)}u,\quad
 \mathrm{Ex}(b)=\frac{e^b-1}b,\quad \mathrm{Ex}_1=\mathrm{Ex}',\quad
 \mathsf N_1(u)=\frac{1-\mathsf L(u)}u
\]
by continuity at zero, and set
\begin{gather*}
 g=\kappa+A\varepsilon z,\quad h=g\varphi+\varepsilon z,\quad
 E=\exp(z\mathsf L(\varepsilon^2z)),\quad
 T=\exp[-(\kappa/\varepsilon)\mathsf L(\kappa\varepsilon)],\\
 \widehat\Lambda=gz/\varphi,\qquad
 \lambda=\varepsilon^3\widehat\Lambda,\qquad
 \Lambda=M\lambda/K=z+\frac{\varepsilon z^2}{\kappa\varphi}.
\end{gather*}
Then $Q^\varepsilon=K^{-3}F=E\widetilde P_2-\widetilde P_1+(T/\varepsilon)W/\kappa^3$, where
\begin{align*}
 \widetilde P_2&=(\varphi g^3-\varepsilon A^2z^2+A^2z^2g)/\kappa^3,\\
 \widetilde P_1&=[\kappa(g^2\varphi^2-\varepsilon^2z^2)+(\kappa+\varepsilon)zgh]/\kappa^3,\\
 W&=-A^2(g^2\varphi^2-\varepsilon^2z^2)
       +[A\varepsilon^4z^3+A^2\varphi g^2(z+\varphi)]/E.
\end{align*}
Thus $Q^0=e^z(1+z^2/\kappa^2)-(1+z)$. These identities are checked symbolically
for indeterminate $M$, rather than by fitting finitely many matrices (checks F1--F4).

\textbf{Exponential remainder bound.}
On $\varepsilon\in[e_-,e_+]$, $\kappa\in[k_-,k_+]$, the quantity $T/\varepsilon^r$
is bounded by
\[
 \frac{\exp(-k_-/e_++k_+^2/2)}{e_+^r}\quad(e_+<k_-/r),
 \qquad \frac{T(e_+,k_-)}{e_-^r}\quad(e_->0).
\]
Use the smaller valid bound, with the first condition interpreted without division by zero
when $r=0$. The first follows from $\log(1+x)\ge x-x^2/2$ and monotonicity of
$\exp(-c/e)/e^r$. For the second, $T$ increases with $e$ and decreases with $k$, since
$x/(1+x)-2\log(1+x)<0$. Moreover
$\partial_\kappa(T/\varepsilon)=-T/(\varepsilon^2A)$ and
$|\partial_\varepsilon(T/\varepsilon)|\le(3\kappa+\varepsilon)T/\varepsilon^3$.
All remainder bounds extend by zero at $\varepsilon=0$. The functions above and their derivatives
are enclosed by convergent power series with explicit geometric or factorial tails.

\subsection{AM2. Uniform Hopf boxes and their identification}
For $X=(\kappa,x,y)$ define
$G(X;\varepsilon)=(\operatorname{Re}Q^\varepsilon,\operatorname{Im}Q^\varepsilon,
\operatorname{Re}\Lambda)$, $z=x+iy$. Let $m$ and $e_m$ be the midpoints of $X$ and a
parameter interval $E_0$, and let $Y$ be an approximate inverse of the midpoint Jacobian.
The parametric Krawczyk enclosure is
\[
 \mathcal K=m-YG(m;e_m)-(Y\partial_\varepsilon G(m;E_0))(E_0-e_m)
                   +(I-YD_XG(X;E_0))(X-m).
\]
If $\mathcal K\subset\operatorname{int}X$, every parameter in $E_0$ has a unique zero in $X$.
Indeed, apply the fundamental theorem of calculus to $YG(m;e)$:
\[
 YG(m;e)-YG(m;e_m)=\int_{e_m}^{e}Y\partial_\varepsilon G(m;t)\,dt.
\]
This justifies forming $Y\partial_\varepsilon G$ before interval multiplication by $e-e_m$.
The ordinary Krawczyk theorem then applies for each fixed $e$. Strict inclusion also bounds
the weighted row sums of $I-YD_XG$ below one, ensuring uniqueness and nonsingularity.

On every stored box we also certify
\[
 0\notin Q_z,\quad |AE+\varepsilon^2z|>0,\quad y>0,\quad
 \operatorname{Im}\Lambda>0,\quad
 \operatorname{Re}(\Lambda_\kappa-\Lambda_zQ_\kappa/Q_z)>0.
\]
These give simplicity, positive frequency, and transversality. All derivatives are obtained by
forward differentiation of the explicit formulas, including the exponential remainders.

The first interval begins at zero. There, the root lies in
$(2j-1)\pi<y<(2j-\tfrac12)\pi$ and has $\kappa>0$, identifying the $j$th root
$\tan y=y$, $\cos y<0$. At every shared parameter endpoint a separate narrow Krawczyk
enclosure is contained in both neighbouring boxes. Therefore their zeros coincide, and
the entire chain is the continuation of this limiting root. This last step is checked by
\texttt{check\_branch\_links.py}; overlap of two large boxes alone would not suffice.

\subsection{AM3. Explicit localization of the spectrum}
For $w=\varphi-1$, $a=|w|$, $u=\operatorname{Re}w$, one has
\[
 \operatorname{Re}\lambda=Au-\frac{u+a^2}{1+2u+a^2}.
\]
Choose the root $\varphi$ of modulus at least $A^{-1/2}$. If $\operatorname{Re}\lambda\ge0$,
then $|\varphi|\ge1$ and $\operatorname{Re}\varphi\ge1$: writing $\varphi=re^{i\theta}$,
$\operatorname{Re}\lambda=\cos\theta(Ar+r^{-1})-(A+1)$; the range
$A^{-1/2}\le r<1$ makes this negative, and for $r\ge1$ it implies
$r\cos\theta\ge(A+1)r^2/(Ar^2+1)\ge1$.
Consequently $u\ge u_{\min}(a)$, the positive root of
$2Au^2+(K+Aa^2)u-a^2=0$, increasing in $a$ and decreasing in $K$.

Here are sufficient inequalities, with explicit constants, to exclude every root with
$a\ge R_0/M$. For $a\le a_1=0.1$, put $\mathcal T=MA^2a^2|K+Aw|$.
The triangle inequality in $F$, divided by $\mathcal T$, excludes a zero if
$\rho_a+\cdots+\rho_f<1$, where
\begin{align*}
 \rho_a&=(1+a_1)(\kappa/(AR_0)+\varepsilon)^2,& \rho_b&=\varepsilon/\kappa,\\
 \rho_c&=(1+a_1)^2[\kappa^2/(A^2R_0^2)+\kappa\varepsilon/(AR_0)],&
 \rho_d&=\varepsilon^2/A^2,\\
 \rho_e&=\frac{\kappa^2+\kappa\varepsilon}{A^2R_0}
              +\frac{(\kappa\varepsilon+\varepsilon^2)(2+a_1)}A,&
 \rho_f&=\frac{N_0T\varepsilon^{-3}+N_1T\varepsilon^{-5}}{A^2R_0^2\kappa},
\end{align*}
with
\begin{align*}
 N_0&=A^2[(K+Aa_1)^2(1+a_1)^2+a_1^2]+Aa_1^3
                 +A^2(1+a_1)^2(K+Aa_1)^2,\\
 N_1&=A^2(1+a_1)(K+Aa_1)^2a_1.
\end{align*}
For example, $|\varphi||K+Aw|^3/\mathcal T\le\rho_a$ follows from
$a\ge R_0/M$; the $(KM+1)S_0$ term gives $\rho_e$. The remaining terms follow by
$|K+Aw|\ge\max(K,Aa)$, $|\varphi|\le1+a$, and $|(A\varphi)^{-M}|\le T$.

For larger $a$, subdivide into $[a_-,a_+]$ and verify
\[
 (1+2u_{\min}(a_-)+a_-^2)^{M_{\rm lo}/2}\underline{|P_2|/M}
    >\overline{|P_1|/M}+\overline{T(|R_1|+|R_2|)/M},\qquad M_{\rm lo}=e_+^{-2}.
\]
The explicit bounds used are
\begin{align*}
 |P_2|&\ge MA^2a^2|K+Aw|-|\varphi||K+Aw|^3-A^2a^2,\\
 |P_1|&\le K(|K+Aw|^2|\varphi|^2+a^2)+(KM+1)a|K+Aw|(K+2Aa+Aa^2),\\
 T(|R_1|+|R_2|)&\le T\{A^2(|K+Aw|^2|\varphi|^2+a^2)+Aa^3
                       +A^2|\varphi||K+Aw|^2(Ma+|\varphi|)\}.
\end{align*}
When $e_->0$, these bounds can also cover $[R_0e_-^2,a_1]$ instead of the six-ratio test.

For branch 1, $K<1$ on all boxes. The first Gershgorin disc lies strictly in the left
half-plane because its centre-to-radius margin is $M(1-K)+6K>0$. The other discs give
$|\lambda|\le2A+2$. Since $|\lambda|\ge Aa^2/(1+a)$, it suffices to check
$a\le a_{\max}=2+2\sqrt2$.
For branches 2 and 3 no restriction on $K$ is needed. With $r=|\varphi|$ one may replace
the lower bound for $|P_2|/M$ by the maximum of that bound and
\[
 A^3r^3-\frac{3A^2r^3+(3A+A^2)r^2}{M}-A^2(1+2A)r^2-A^2(2+A)r
       -\frac{(1+2A^2)r}{M}-A^2(1+M^{-1}).
\]
This follows by expanding $P_2$ and using $|\varphi+M|\ge M+1$.
The finite subdivision ends at $a_{\rm big}=12$. For $r\ge r_b=a_{\rm big}-1$, set
$\alpha=K/2+3A/2$, $\beta=K/4+15A/4$, and $d=M^{-1}$. Bounds for the tail are
\begin{align*}
 C_1&=K(\alpha^2+9/16)d+\tfrac32(K+d)\alpha\beta,\\
 C_2&=T\{[A^2(\alpha^2+9/16)+27A/16+A^2\alpha^2]d
                        +\tfrac32A^2\alpha^2\},\\
 c&=A^3-3A^2d-(3A+A^2)d/r_b-A^2(1+2A)/r_b\\
  &\hspace{1em}-A^2(2+A)/r_b^2-(1+2A^2)d/r_b^2-A^2(1+d)/r_b^3.
\end{align*}
Take lower endpoints for positive terms in $c$ and upper endpoints for subtracted terms.
Then $|P_1|/M\le C_1r^4$, $T(|R_1|+|R_2|)/M\le C_2r^4$, and $|P_2|/M\ge cr^3$.
The certified inequality $r_b^{M_{\rm lo}-1}c>C_1+C_2$ excludes the entire tail.
Thus all possible unstable or imaginary eigenvalues are localized to $|z|<R_0$.

\subsection{AM4. Counting zeros without floating-point angles}
The following argument-principle calculation uses balls throughout. On a contour edge let
$B_e$ enclose $Q^\varepsilon$ for the whole edge and parameter box. At the midpoint
parameters let $v_0,v_1$ enclose the endpoint values. Subdivide until
$\operatorname{Re}(B_e/v_0)>0$. The actual edge image, rotated by its initial value,
lies in the right half-plane; its argument increment is therefore
$\arg(v_1/v_0)$, computed in Arb. The sum divided by $2\pi$ must contain a unique
integer. Nonvanishing on all edges and homotopy invariance make this count valid
throughout the parameter box. The former floating-point sum of angles is not used.

For branch 1 take $R_0=8$, $\sigma=0.3$, and a convex polygon enclosing
$D=\{|z|\le R_0,\ \operatorname{Re}z\ge-\sigma\}$.
The polygon's supporting half-planes are checked in ball arithmetic, including the
slight outward displacement of its left chord. If $x\le-\sigma$ and $|z|\le R_0$,
\[
 \operatorname{Re}\Lambda\le
 \max(-\sigma+\varepsilon\sigma^2/\kappa,-R_0+\varepsilon R_0^2/\kappa)
 +\frac{\varepsilon^3R_0^3}{\kappa(1-\varepsilon^2R_0)}<0.
\]
The bound follows by expanding $1/\varphi$ about one and maximizing the resulting convex
quadratic in $x$. Each box has winding count four. Two zeros are $z=0$, counted at least
twice because $F(1)=F'(1)=0$; the other two are the certified Hopf pair. Localization
and the displayed inequality leave no other eigenvalue with nonnegative real part on $H$.

For branches 2 and 3 use polygons following level curves of $\operatorname{Re}\Lambda$.
The constants are respectively
\[
 (R_0,Y,c_L,c_R)=(20,20.4,-0.3,0.25),\qquad(26,26.6,-0.15,0.25).
\]
Piecewise linear left and right graphs $\ell(y),r(y)$ have strictly negative and positive
$\operatorname{Re}\Lambda$, respectively, certified edge by edge. They lie in
$\mathcal R=[x_{\min},R_0]\times[-Y,Y]$. On this rectangle, with
$\rho=\max_{\mathcal R}|z|$, the check
\[
 \frac{\varepsilon}{\kappa}\frac{\rho(2+\varepsilon^2\rho)}{(1-\varepsilon^2\rho)^2}<1
\]
implies $\partial_x\operatorname{Re}\Lambda>0$. To exclude the part of $|z|\le R_0$
left of $x_{\min}=-\sigma$, put $q=\varepsilon^2R_0\le2/5$ and
$g_0(x)=x+\varepsilon x^2/\kappa+\varepsilon^3|x|^3/[\kappa(1-q)]$.
The sharper bound is $\operatorname{Re}\Lambda\le\max(g_0(-\sigma),g_0(-R_0))<0$.
Indeed the upper bound at fixed $x$, as a function of $t=y^2$, is
$x+(\varepsilon/\kappa)(x^2-t)+\varepsilon^3(x^2+t)^{3/2}/[\kappa(1-q)]$;
its derivative is nonpositive when $q\le2/5$, and $g_0$ is convex for $x<0$.

The polygons $P=\{\ell(y)<x<R_0,|y|<Y\}$ and
$P_-=\{\ell(y)<x<r(y),|y|<Y\}$ have counts $2j+2$ and four. Localization and
monotonicity imply that all nonnegative spectrum is in $P$, and all imaginary spectrum
is in $P_-$. Its four zeros are exactly $0,0,z_H,\bar z_H$.
The remaining $2j-2$ zeros in $P\setminus P_-$ give eigenvalues with positive real part;
the real exceptional points of $F$ lie outside this region. This proves H2 and instability.
No claim about simplicity of these additional zeros is needed.

\subsection{AM5. A regular formula for the nonlinear coefficient}
Let $q_1=1$, $pB=i\omega p$, $pq=1$, $r(\mu)=(\mu-B)^{-1}Sq$, where $r(0)\in H$,
and $\sigma(v)=pSv$. The quadratic projection formula of Section 3 gives exactly
\begin{align*}
 \ell_1=\omega^{-1}\operatorname{Re}\big[&\sigma(r(0))+\sigma(\overline{r(0)})
 +2\operatorname{Re}(r(0)_1)\sigma(q)\\
 &+\sigma(r(2i\omega))+r(2i\omega)_1\sigma(\bar q)\big].
\end{align*}
For example $w_{11}=-(r(0)+\overline{r(0)})$, $w_{20}=2r(2i\omega)$, which proves
this identity by substitution. In computing a resolvent, the bulk rows, last row and
$\sum_l l r_l=0$ replace the first row: multiplying by the conserved mass row proves
the latter automatically for $\mu\ne0$, and selects the restricted inverse for $\mu=0$.

Here are all the blocks needed to evaluate the formula without singular cancellation.
Write $T_q=T/\varepsilon$ and
\begin{align*}
 C_1&=\frac{g-A^2\varphi T_q/E}{gE/\varphi+\varepsilon^2zAT_q/E},&
 \widehat C_2&=\frac{AE+\varepsilon^2z}{gE+\varepsilon^2zA\varphi T_q/E},\\
 q_l&=C_1E\varphi^{-l}+\varepsilon^{-1}\widehat C_2(A\varphi)^{l-M}.&&
\end{align*}
For the unnormalized adjoint put $\psi=(A\varphi)^{-1}$,
$\theta=\widehat\Lambda/(\kappa+\varepsilon^2\widehat\Lambda)$,
$c_A=(\kappa+\varepsilon)\theta+\kappa$, and
\begin{align*}
 a_1&=c_A\psi+\varepsilon^2T_q(Az-\varepsilon\widehat\Lambda)/E,\\
 a_2&=-c_A\varphi+E(g/\varphi+\varepsilon^2\widehat\Lambda),\\
 p_l&=a_1\varphi^l+a_2\psi^l-(1-\varepsilon^2\theta)(a_1\varphi+a_2\psi).
\end{align*}
Divide by $N_p=\sum_l p_lq_l$ when forming $\sigma$. For $l<M$,
\[
 (pS)_l=a_1\varepsilon^2z\varphi^l-a_2\frac{\varepsilon g}{A\varphi}\psi^l
             -\varepsilon^2\theta(a_1\varphi+a_2\psi),\qquad(pS)_M=0.
\]
These formulas follow by substitution in the second-order bulk recurrence and its
boundary rows; the nonsingular normalization is checked by the ball evaluation.

For $\mu=\varepsilon^3m$, where $m=0$ or $2\widehat\Lambda$, put
\[
 c_m=\kappa+\varepsilon^2m,\quad
 \zeta=\frac{-2m}{c_m+\sqrt{c_m^2+4\varepsilon m}},\quad
 s_1=1+\varepsilon^2\zeta,\quad s_2=A/s_1,
\]
and $\Xi=(\kappa-\varepsilon\zeta)/(1+\varepsilon^2\zeta)+A\varepsilon z-\varepsilon^2m$.
Then
\begin{align*}
 r_l={}&\frac{zC_1E}{\varepsilon(m-\widehat\Lambda)}\varphi^{-l}
 +\frac{\widehat C_2g}{\varepsilon\Xi}
       \frac{(A\varphi)^{l-M}-s_2^{l-M}}{A\varphi-s_2}\\
 &+\frac{\widehat h_1}{\varepsilon}s_1^l
       +\frac{\widehat H_2}{\varepsilon^2}s_2^{l-M}.
\end{align*}
The two coefficients are obtained from the last row and the mass condition, scaled
by $\varepsilon$ and $\varepsilon^5$ respectively. These are explicit two-by-two linear
equations. The divided difference follows from
$(w-s_2)(w+s_2-A-1-\mu)=-(\mu-\lambda)w$, $w=A\varphi$;
it removes the subtraction of two individually larger boundary layers.

For completeness the finite sums used in these blocks are evaluated as follows.
For $\beta=1+\varepsilon^2b_0$, $b=M\log\beta=b_0\mathsf L(\varepsilon^2b_0)$,
and $\gamma=A(1+\varepsilon^2g_0)$, geometric summation gives
\begin{align*}
 \sum_{l=1}^{M-1}(\beta_1\beta_2)^l
    &=\varepsilon^{-2}[\mathrm{Ex}(b_1+b_2)\mathsf L(\beta_1\beta_2-1)-\varepsilon^2],\\
 \sum_{l=1}^{M-1}\beta^l\gamma^{l-M}
    &=\frac{e^b-\beta\gamma^{1-M}}{\beta\gamma-1},\\
 \sum_{l=1}^{M}l\beta^l
    &=\varepsilon^{-4}\beta[\mathsf L(u)^2\mathrm{Ex}_1(b)
                              +\varepsilon^2e^b\mathsf N_1(u)],\quad u=\beta-1,\\
 \sum_{l=1}^{M}l\gamma^{l-M}
    &=\varepsilon^{-3}(\gamma/d_\gamma-\varepsilon\gamma/d_\gamma^2)
                              +\gamma^{1-M}/(\gamma-1)^2,
 \quad d_\gamma=(\gamma-1)/\varepsilon.
\end{align*}
Divided differences of these identities give the corresponding sums with
$(\gamma_1^{l-M}-\gamma_2^{l-M})/(\gamma_1-\gamma_2)$; their regular parts are
$-\beta e^b/(\varepsilon^2d_1d_2)$ for the pair sum, where
$d_i=(\beta\gamma_i-1)/\varepsilon$, and
\[
 \varepsilon^{-4}\left[-\frac{1-\varepsilon^2}{d_1d_2}
       +\frac{\varepsilon(d_1+d_2)}{d_1^2d_2^2}\right]
\]
for the mass sum, now with $d_i=(\gamma_i-1)/\varepsilon$.
The remainders are bounded by the supremum of the derivative of the exponentially small
part along the segment $[\gamma_1,\gamma_2]$. Products with the layer
$\psi'^l$, $\psi'=(1+\varepsilon^2p_0)/A$, use
$\sum_{l=1}^{M-1}(\psi'\beta)^l=\psi'\beta[1-(\psi'\beta)^{M-1}]/(1-\psi'\beta)$;
products with a layer at $M$ are exponentially small. The bound in AM1 controls all
remainders, including coincident bases. Direct summation and dense linear solves check
these identities entrywise (F5--F6), without using the earlier closed-form implementation.

Track a coefficient as $v\varepsilon^p$, with integer $p$ and regular $v$.
The five terms in the bracket above have order $\varepsilon$; all negative powers cancel
before interval evaluation. Thus
\[
 \widehat L=\varepsilon^2\ell_1(q_1=1)
      =\frac{\operatorname{Re}(\varepsilon^{-1}\text{bracket})}
                   {\operatorname{Im}\widehat\Lambda}
\]
is continuous at $\varepsilon=0$. To control interval overestimation, regard each
exponential remainder as an independent variable $\nu_i$ in a disk valid on the whole box.
For the resulting regular expression $\Phi$, at any point $\xi$ of the box with midpoint $c$,
\[
 \Phi(\xi;\nu(\xi))\in\Phi(c;\mathcal D)
              +D_\xi\Phi(\text{box};\mathcal D)(\xi-c).
\]
This is the mean-value formula at fixed $\nu(\xi)$, not differentiation of a numerical
remainder bound. The same disks are recorded and replayed in the box and midpoint evaluations.

\subsection{AM6. Finite certificates covering the infinite range}
All computations use python-flint 0.9 at 192 bits, except the higher precision limiting
evaluation. The exact rational endpoint stored in the certificates is
\[
 e_*=0.078811099999999995>161^{-1/2}.
\]
\begin{center}
\begin{tabular}{c r c c c}
\toprule
Branch & Boxes & $\widehat L$ & Spectral counts & Integer range\\
\midrule
1 & 562 & $<0$ & $4$ & $M\ge161$\\
2 & 589 & $<0$ & $N(P)=6,\ N(P_-)=4$ & $M\ge161$\\
3 & 834 & $<0$ & $N(P)=8,\ N(P_-)=4$ & $M\ge161$\\
\bottomrule
\end{tabular}
\end{center}
Each line is re-proved from its midpoint and radius, including the Krawczyk inclusion,
all nondegeneracy conditions, and the coefficient sign. No integer sampling is used for
this infinite range. Exact rational endpoint checks prove coverage, and AM2 identifies
the whole branch. The first box $[0,10^{-5}]$ covers every $M\ge10^{10}$.
For $M\le160$, the direct reduced-matrix certificate from Section 11 is rerun at every
integer: 146, 117 and 74 points, respectively, including its rigorous eigensolver.
This proves Theorem D's coefficient and spectral assertions for every relevant integer.

At $\varepsilon=0$, symbolic evaluation of the same blocks gives
\[
 \lim_{M\to\infty}\frac{\ell_1(q_1=1)}M
 =-\frac{R(R-1)(7R^2-34R+23)}{t(R+1)(5R-3)((R-1)^3+R+1)},\quad
 R=\sqrt{1+t^2},\quad t=t_j.
\]
To check the identity exactly, substitute
$R=(4\tau^4+1)/(4\tau^2)$, $t=(4\tau^4-1)/(4\tau^2)$,
$\kappa=(2\tau^2-1)/(2\tau)$, $e^{it}=-(1+it)/R$ in the blocks.
Clearing denominators reduces the difference to the zero polynomial
(\texttt{limit\_symbolic.py}). Real-root isolation checks that the normalization and
resolvent divisors, including their own denominators, never vanish for $\tau\ge2$.
All the Hopf roots have $\tau>2.12$. Since $R\ge65/16>(17+8\sqrt2)/7$,
the displayed coefficient is negative. In particular the first three limits are
$-0.00931139838\ldots$, $-0.00794101725\ldots$, $-0.00377769679\ldots$.
This is a limit of the discrete-system coefficient, not a deduction from the formal renewal model.

\subsection{AM7. Exact onset and the finite-to-infinite junction}
The second chart makes the limit $K\to\infty$ regular. Put $s=1/A$, allow real $M>3$,
and use $\varphi^M=\exp(M\log\varphi)$ in the right half-plane. Define
\begin{align*}
 f_1&=\varphi^M+s(\varphi-1),&f_2&=s-\varphi+s^{M-1}\varphi^{1-M},\\
 G_1&=\varphi(s-\varphi)[(1-s)\varphi+(M(1-s)+s)(\varphi-1)],\\
 G_2&=(s-\varphi)[M(\varphi-1)+(1-s)\varphi+\varphi(\varphi-1)],\\
 H_1&=(\varphi-1)[s(1-s)+(M(1-s)+s)(s-\varphi)],\\
 H_2&=(\varphi-1)[(s+M\varphi)(s-\varphi)+s(1-s)\varphi].
\end{align*}
Then $F/A^3=f_1(H_1-s^{M-1}\varphi^{-M}H_2)-(G_1-\varphi^MG_2)f_2$ and
$\lambda/A=(1-s/\varphi)(\varphi-1)$. These follow by substitution in the
boundary determinant defining $F$. At $s=0$ the Hopf equations reduce to $u=0$ and
\[
 G_\infty(1+iv;M)=0,\qquad
 G_\infty(\varphi;M)=-\varphi^{M-2}[\varphi^3+M(\varphi-1)^2]+(M+1)\varphi-M.
\]
The factors $s^{M-1}$ are $C^1$ at zero. For $0\le s\le b<1$ and $p=M-1>2$,
their values and derivatives are bounded by $b^{p_-}$, $p_+b^{p_--1}$ and
$b^{p_-}|\log b|$, the last provided $b\le e^{-1/p_-}$.

The script \texttt{check\_onsets.py} reuses the explicit onset chart from the companion
folder, without assuming any result about higher branches. It first encloses
$(M,u,v)$ as functions of $s\in[0,0.0002]$, with $dM/ds>0$ and $v>0$.
It then encloses $(s,u,v)$ as functions of $M$ up to $162$, with $0<s<1$ and $v>0$.
Every stored line and every consecutive endpoint is re-proved. A separately enclosed
point at the chart transition belongs to both unique-zero boxes; a final sub-box maps
inside the last $\varepsilon$-chart box under
\[
 \varepsilon=M^{-1/2},\quad\kappa=(s^{-1}-1)\sqrt M,\quad z=M(\varphi-1).
\]
Every integer Hopf point used for $M\le161$ is also enclosed afresh and contained in
this same curve. The first $s$-boxes lie in the intervals for $M_j^*$ in Theorem D.
Their positive slope, the graph over $M$, and the nonsingular local Hopf systems show
that this arc is the entire connected component: there is no finite interior point
where an additional arc can attach. Its two ends leave the domain at $K=\infty$ and
$M=\infty$. This proves the exact onsets, including the previously uncertified $87$.

Finally, for $3\le M\le14$, \texttt{check\_small\_stability.py} derives the reduced
Jacobian directly from the vector field, computes its characteristic polynomial in
$\mathbb Z[K][\lambda]$, and verifies positivity on $K>0$ of every leading Hurwitz
determinant. Positive polynomial coefficients or exact real-root isolation establish
these inequalities. The Hurwitz criterion proves stability throughout this parameter range.

\subsection{AM8. Consequences, verification, and scope}
The Hopf theorem now applies at every integer point of the three curves. The sign
and crossing direction give the asserted side and stability on the centre manifold.
For branch 1 all transverse eigenvalues are stable. For $j=2,3$, the monodromy matrix
of the small orbit tends to $\exp(2\pi B|_H/\omega)$; a transverse positive-real-part
eigenvalue therefore gives a Floquet multiplier outside the unit circle. Positivity of
the equilibrium ensures that sufficiently small cycles have positive concentrations.

The complete recheck is \texttt{all-M/scripts/recheck\_all.py}; its logs, generated onset
certificates and direct matrix certificates are in
\texttt{all-M/verification/recheck-2026-10-25/}. The detailed audit is
\texttt{all-M/verification.md}. Corrections made during the audit include the fully
interval winding count, certified contour containment, and explicit branch junctions.
The proof uses correctness of Arb, exact polynomial arithmetic and the standard
finite-dimensional Hopf and Krawczyk theorems. It makes no assertion about cycles far
from bifurcation, a global stability threshold in $\kappa$, or the unresolved higher
branches. Those issues are not needed for the complete local answer on branches 1--3.
